\section*{Hoofdstuk \ref{ch:b2:angiosperm-reproduction} --- Geslachtelijke voortplanting van bloemplanten}

\begin{solution}{exo:b2:angiosperm-reproduction:1}
Kelkbladen (kelk), kroonbladen (kroon), meeldraden (androecium),
\omterm{def:b2:angiosperm-reproduction:flower}{vruchtbladen} (gynoecium). \omterm{def:b2:angiosperm-reproduction:flower}{Meeldraad}: helmdraad en helmknop. \omterm{def:b2:angiosperm-reproduction:flower}{Vruchtblad}:
\omterm{def:b2:angiosperm-reproduction:flower}{stempel}, stijl, \omterm{def:b2:angiosperm-reproduction:flower}{vruchtbeginsel} met \omterm{def:b2:plant-life-cycles:heterospory}{zaadbeginsels}. Al deze delen zijn
diploïde \omterm{def:b2:plant-life-cycles:cycles}{sporofyt}; alleen de \omterm{prop:b2:angiosperm-reproduction:pollen}{stuifmeelkorrel} (in de helmknop) en de
\omterm{prop:b2:angiosperm-reproduction:embryosac}{embryozak} (in het \omterm{def:b2:plant-life-cycles:heterospory}{zaadbeginsel}) zijn haploïde \omterm{def:b2:plant-life-cycles:cycles}{gametofyt}.
\end{solution}

\begin{solution}{exo:b2:angiosperm-reproduction:2}
\omterm{def:b2:plant-life-cycles:heterospory}{Microspore} $n$; vegetatieve cel $n$; spermacel $n$; eicel $n$; poolkern
$n$; zygote $2n$; endosperm $3n$; \omterm{def:b2:angiosperm-reproduction:seed}{zaadhuid} $2n$ (maternale
integumenten); vruchtwand $2n$ (maternaal \omterm{def:b2:angiosperm-reproduction:flower}{vruchtbeginsel}).
\end{solution}

\begin{solution}{exo:b2:angiosperm-reproduction:3}
Wind (grassen, eik, berk, hazelaar, weegbree): kleine groene \omterm{def:b2:angiosperm-reproduction:flower}{bloemen},
zonder kroonbladen, geur of \omterm{def:b2:angiosperm-reproduction:pollination}{nectar}, helmknoppen die aan lange helmdraden
bungelen, veervormige \omterm{def:b2:angiosperm-reproduction:flower}{stempels}, enorme hoeveelheden glad droog \omterm{def:b2:plant-life-cycles:heterospory}{stuifmeel},
bloei vóór het blad. Bij (klaver, salie, vingerhoedskruid, lavendel,
appel): gekleurde kroonbladen, vaak met ultraviolette nectarwijzers, een
landingsplatform, geur, \omterm{def:b2:angiosperm-reproduction:pollination}{nectar}, kleverig gebeeldhouwd \omterm{def:b2:plant-life-cycles:heterospory}{stuifmeel} in matige
hoeveelheden, bloei bij daglicht.
\end{solution}

\begin{solution}{exo:b2:angiosperm-reproduction:4}
De twee spermacellen van één \omterm{thm:b2:angiosperm-reproduction:tube}{stuifmeelbuis} versmelten, de ene met de eicel
(zygote, $2n$, het embryo) en de andere met de centrale cel (endospermkern,
$3n$, de voedselvoorraad). Het endosperm begint pas na de bevruchting te
groeien, zodat de plant alleen \omterm{def:b2:plant-life-cycles:heterospory}{zaadbeginsels} bevoorraadt die een embryo
bevatten, in tegenstelling tot de naaktzadige, die zijn voedselvoorraad
vooraf opbouwt.
\end{solution}

\begin{solution}{exo:b2:angiosperm-reproduction:5}
$25/1.2 = \qty{21}{h}$. Draden die op dag 5, 6 en 7 verschijnen, treffen
\omterm{def:b2:plant-life-cycles:heterospory}{stuifmeel} aan; die van dag 8 tot 12 niet: $3/8$ van de draden wordt
bestoven, en de kolf is voor vijf achtste leeg.
\end{solution}

\begin{solution}{exo:b2:angiosperm-reproduction:6}
$S_1S_2\times S_1S_2$: 0. $S_1S_2\times S_1S_3$: $\tfrac12$ (alleen
\omterm{def:b2:plant-life-cycles:heterospory}{stuifmeel} $S_3$ groeit). $S_1S_2\times S_3S_4$: alles. Nakomelingen van de
tweede kruising: eicellen $S_1$ of $S_2$, spermacel $S_3$: $S_1S_3$ en
$S_2S_3$, elk de helft.
\end{solution}

\begin{solution}{exo:b2:angiosperm-reproduction:7}
Een plant draagt 2 van de $n$ \omterm{def:b2:meiosis-heredity:vocabulary}{allelen} en wijst \omterm{def:b2:plant-life-cycles:heterospory}{stuifmeel} af dat een van
beide draagt: compatibele fractie $(n - 2)/n$. Een nieuw \omterm{def:b2:meiosis-heredity:vocabulary}{allel} wordt door
geen enkele stamper afgewezen, zodat zijn \omterm{def:b2:plant-life-cycles:heterospory}{stuifmeel} meer zaden verwekt dan
gemiddeld en het zich verspreidt tot het even gewoon is als de rest;
hetzelfde voordeel beschermt elk zeldzaam \omterm{def:b2:meiosis-heredity:vocabulary}{allel} tegen verlies. Selectie
stapelt dus \omterm{def:b2:meiosis-heredity:vocabulary}{allelen} op, en natuurlijke populaties van zelfincompatibele
soorten dragen er tientallen.
\end{solution}

\begin{solution}{exo:b2:angiosperm-reproduction:8}
De twee \omterm{prop:b2:angiosperm-reproduction:embryosac}{poolkernen} zijn kopieën van één \omterm{def:b2:plant-life-cycles:heterospory}{megaspore}, zodat de centrale cel
$YY$ of $yy$ is (elk de helft); de spermacel is $Y$ of $y$ (elk de
helft): endosperm $YYY$, $YYy$, $Yyy$, $yyy$ met elk $\tfrac14$; drie
kwart geel. Bestoven door $yy$: $YYy$ en $yyy$, elk de helft --- de helft
van de korrels geel.
\end{solution}

\begin{solution}{exo:b2:angiosperm-reproduction:9}
Gerstkorrels worden gehalveerd; de helften met embryo verteren hun
zetmeel, de helften zonder embryo niet, tenzij er gibberelline wordt
toegevoegd, waarna ze het wel doen. (a) Het embryo is de bron van het
signaal, niet van het enzym. (b) Gibberelline is het signaal, dat op
zichzelf volstaat. (c) Het amylase wordt gemaakt door de aleuronlaag van
het endosperm, die op het hormoon reageert. De helft zonder embryo scheidt
signaal en respons: zonder haar kon men niet weten of het embryo het
zetmeel zelf verteert.
\end{solution}

\begin{solution}{exo:b2:angiosperm-reproduction:10}
Wind: $10^{6}\times\qty{1}{\micro g} = \qty{1}{g}$ per \omterm{def:b2:plant-life-cycles:heterospory}{zaadbeginsel}. Dier:
\qty{1}{mg} \omterm{def:b2:plant-life-cycles:heterospory}{stuifmeel} plus \qty{0.5}{mg} suiker $= \qty{1.5}{mg}$
per \omterm{def:b2:plant-life-cycles:heterospory}{zaadbeginsel} --- ongeveer zevenhonderd keer goedkoper. Windbestuiving
heeft geen partner nodig: ze werkt in het vroege voorjaar voordat insecten
vliegen, op koude, winderige en open plaatsen, 's nachts en in de regen,
in dichte bestanden van één soort, en ze kan niet worden bedrogen door
nectardieven of in de steek gelaten door een bestuiver die uitsterft.
\end{solution}

\begin{solution}{exo:b2:angiosperm-reproduction:11}
Zaden die aan de voet van de ouderboom vallen, landen in zijn schaduw,
tussen zijn wortels, en waar zijn specifieke zaadeters en ziekteverwekkers
zich concentreren; bijna alle sterven. Een vogel draagt een \omterm{def:b2:angiosperm-reproduction:seed}{zaad} kilometers
ver (\qty{12}{km} in twintig minuten) en laat het, met mest, in een heg of
een open plek vallen; zelfs als maar één \omterm{def:b2:angiosperm-reproduction:seed}{zaad} op honderd op een goede plek
landt, is dat meer dan alle zaden onder de boom samen, en de populatie
verspreidt zich met de snelheid van vogels, niet van vallend fruit. Het
vruchtvlees is de prijs van het vervoersbewijs.
\end{solution}

\begin{solution}{exo:b2:angiosperm-reproduction:12}
De insluiting laat het \omterm{def:b2:plant-life-cycles:heterospory}{stuifmeel} door maternaal weefsel groeien, zodat
de stamper het kan testen (zelfincompatibiliteit, afwijzing van vreemde
soorten), veel buizen met elkaar kan laten wedijveren zodat de snelste het
\omterm{def:b2:angiosperm-reproduction:seed}{zaad} verwekt, de \omterm{def:b2:plant-life-cycles:heterospory}{zaadbeginsels} kan beschermen tegen uitdroging en
planteneters, en het \omterm{def:b2:angiosperm-reproduction:flower}{vruchtbeginsel} kan omvormen tot een \omterm{def:b2:angiosperm-reproduction:seed}{vrucht} die dieren
voor de verspreiding inschakelt; de \omterm{prop:b2:angiosperm-reproduction:double}{dubbele bevruchting} bevoorraadt dan
alleen bevruchte \omterm{def:b2:plant-life-cycles:heterospory}{zaadbeginsels}. De kosten: een \omterm{prop:b2:angiosperm-reproduction:pollen}{stuifmeelkorrel} moet een
\omterm{def:b2:angiosperm-reproduction:flower}{stempel} bereiken in plaats van het \omterm{def:b2:plant-life-cycles:heterospory}{zaadbeginsel} zelf, zodat een stijl, een
buis en de geleiding ervan nodig zijn, en de \omterm{def:b2:angiosperm-reproduction:seed}{vrucht} is een zware
investering. De balans viel zo sterk in het voordeel van de bedektzadigen
uit dat ze in honderd miljoen jaar van niets tot negen tiende van de
plantensoorten groeiden.
\end{solution}

\begin{solution}{pb:b2:angiosperm-reproduction:1}
\textbf{1.} Op A ($S_1S_2$): \omterm{def:b2:plant-life-cycles:heterospory}{stuifmeel} A $f = 0$; B ($S_1$, $S_3$) $f =
\tfrac12$; C ($S_3$, $S_4$) $f = 1$.
\textbf{2.} Op B ($S_1S_3$): A $\tfrac12$, B 0, C $\tfrac12$. Op C
($S_3S_4$): A 1, B $\tfrac12$, C 0.
\textbf{3.} 50 compatibele korrels; $0.1\times 50 = 5$ \omterm{def:b2:plant-life-cycles:heterospory}{zaadbeginsels}.
\textbf{4.} Van C: 100 compatibel, $0.1\times 100 = 10$, d.w.z.\ alle
tien \omterm{def:b2:plant-life-cycles:heterospory}{zaadbeginsels}; van A: geen.
\textbf{5.} \omterm{def:b2:angiosperm-reproduction:seed}{Vruchten} na bezoeken van C (10 zaden) blijven hangen; na
bezoeken van B (5 zaden) net wel; na bezoeken van A worden ze afgestoten.
Een boomgaard met alleen A draagt niets.
\textbf{6.} De helft van de \omterm{def:b2:angiosperm-reproduction:flower}{bloemen} zet \omterm{def:b2:angiosperm-reproduction:seed}{vrucht} (die bezocht vanuit C): 2500
\omterm{def:b2:angiosperm-reproduction:seed}{vruchten} per boom, tien keer de 250 die nodig zijn --- een volle oogst, en
de boom zal het overschot afstoten.
\textbf{7.} Hij geeft gedurende het hele bloeiseizoen van elk ras
compatibel \omterm{def:b2:plant-life-cycles:heterospory}{stuifmeel} af; maar als hij beide $S$-allelen met een ras
deelt, bestuift hij daarop niets.
\textbf{8.} $8\times 10^{7}/(7\times 86400) = 132$ korrels per vierkante
meter per seconde.
\textbf{9.} $132/0.25 = 530$ korrels per kubieke meter.
\textbf{10.} $132\times 10^{-4} = 0.013$ per seconde: één korrel per
\qty{76}{s}.
\textbf{11.} $0.013\times 86400 \approx 1100$ korrels per dag; de eerste
buis die het \omterm{def:b2:plant-life-cycles:heterospory}{zaadbeginsel} bereikt, bevrucht het, en het \omterm{def:b2:plant-life-cycles:heterospory}{zaadbeginsel} is
daarna voor de rest gesloten.
\textbf{12.} \qty{20}{h} na de \omterm{def:b2:angiosperm-reproduction:pollination}{bestuiving}.
\textbf{13.} \omterm{def:b2:plant-life-cycles:heterospory}{Stuifmeel} op dag 0--7, draden vanaf dag 5: alleen de dagen 5,
6 en 7 overlappen --- $3/8$ van de afgifteperiode, en draden die na dag 7
verschijnen, krijgen niets; de kolf blijft grotendeels leeg. Droogte tijdens
de bloei van de draden is de klassieke oorzaak van een mislukte maïsoogst.
\textbf{14.} Embryo $Yy$; endosperm $YYy$.
\textbf{15.} Geel, omdat $Y$ \omterm{def:b2:meiosis-heredity:vocabulary}{dominant} is: het kenmerk van de stuifmeelouder
komt tot uiting in het \omterm{def:b2:angiosperm-reproduction:seed}{zaad} op de moederplant --- xenie.
\textbf{16.} Centrale cel $YY$ of $yy$: endosperm $YYy$ (de helft) en
$yyy$ (de helft).
\textbf{17.} $YYy$: 20 eenheden; $Yyy$: 10 eenheden; een $yyY$ met de $Y$
van de vader is hetzelfde \omterm{def:b2:meiosis-heredity:vocabulary}{genotype} $Yyy$, 10 eenheden. Twee paternale
doses kunnen niet bestaan: één spermacel bevrucht de centrale cel, en het
endosperm bestaat altijd uit twee maternale genomen op één paternaal.
\textbf{18.} Het is de zetmeelvoorraad; $0.3/0.33 = \qty{91}{\%}$ van de
korrel.
\textbf{19.} De voorraad wordt pas na de bevruchting opgebouwd, zodat er
geen middelen naar \omterm{def:b2:plant-life-cycles:heterospory}{zaadbeginsels} gaan die geen embryo dragen; de \omterm{def:b2:plant-life-cycles:cycles}{gametofyt}
van de naaktzadige wordt vooraf opgebouwd en gaat verloren als de
\omterm{def:b2:angiosperm-reproduction:pollination}{bestuiving} mislukt.
\textbf{20.} $0.9\times 600 = 540$ korrels per kolf en per plant;
$4320$ per vierkante meter; $4320\times 0.3 = \qty{1.3}{kg}$ endosperm
per vierkante meter.
\textbf{21.} \qty{13}{t/ha}.
\textbf{22.} Koolstof in het graan $0.5\times 1.2 = \qty{0.6}{kg/m^2}$;
koolstof in het endosperm $0.45\times 1.3 = \qty{0.58}{kg/m^2}$:
consistent (het kleine restant is embryo en \omterm{def:b2:angiosperm-reproduction:seed}{zaadhuid}).
\textbf{23.} $8\times 10^{7}/4320 \approx \num{18500}$ \omterm{prop:b2:angiosperm-reproduction:pollen}{stuifmeelkorrels} per
geoogste maïskorrel.
\textbf{24.} Zijn eigen stuifmeelwolk is dun en waait weg, en zijn pluim
geeft het meeste \omterm{def:b2:plant-life-cycles:heterospory}{stuifmeel} af voordat zijn eigen draden verschijnen, zodat
veel draden nooit worden bestoven en de kolf gaten vertoont.
\textbf{25.} Blok A: de helft van de \omterm{def:b2:angiosperm-reproduction:flower}{bloemen} zet \omterm{def:b2:angiosperm-reproduction:seed}{vrucht}, 2500 per boom, een
volle oogst. Akker: 4320 korrels per vierkante meter, \qty{13}{t/ha}.
\end{solution}
